%\pdfoutput=1
\documentclass[a4paper,11pt]{article}
%\let\ifpdf\relax
%\usepackage{ifpdf}
\usepackage{jheppub}
%\usepackage{natbib}
\usepackage{bm}
\usepackage{pgfplots}
\pgfplotsset{compat=newest}
\usetikzlibrary{shadows.blur}
\usepackage{float}
%\JHEP{00(2002)000}
%\usepackage{graphicx}% use this package if an eps figure is included.
%\usepackage{subcaption}
%\JHEPspecialurl{http://jhep.sissa.it/JOURNAL/JHEP3.tar.gz}

%\usepackage{epsfig,multicol,amsmath,amssymb}
%\usepackage[T1]{fontenc}
%%%%%%%%%%%%%%%%%%%%%%%% latex 2 eps %%%%%%%%%%%%%%%%%%%%%%%%%%
%\usepackage{epstopdf}
%\DeclareGraphicsRule{.tif}{png}{.png}{`convert #1 `basename #1  .tif`.png} 
%\usepackage{bm}  % for the command \bm (boldface greek letters)
%\usepackage{amsmath}     % defines the command \dfrac (fractions in displaystyle
%\usepackage{epsfig,epsf}
%\usepackage{showkeys}
%\usepackage{graphicx}
%\usepackage{graphicx}
%\usepackage{fancybox}
%\usepackage{epsfig,multicol}
%\usepackage{rotating}
%\usepackage{showkeys}
%\textwidth=13.5cm
%\textheight=22cm
%\topmargin=0cm
%\oddsidemargin=0cm
%\usepackage{cite}
%\usepackage{lmodern}
%\usepackage{lipsum}
\voffset1.5cm



%\renewcommand{\thefootnote}{\fnsymbol{footnote}}

\def\beq{\begin{equation}}
\def\eeq{\end{equation}}
\def\bea{\begin{eqnarray}}
\def\eea{\end{eqnarray}}
%

\def\eq#1{{Eq.~(\ref{#1})}}
\def\fig#1{{Fig.~\ref{#1}}}
\newcommand{\bas}{\bar{\alpha}_S}
\newcommand{\as}{\alpha_S}
\newcommand{\bra}[1]{\langle #1 |}
\newcommand{\ket}[1]{|#1\rangle}
\newcommand{\bracket}[2]{\langle #1|#2\rangle}
\newcommand{\intp}[1]{\int \frac{d^4 #1}{(2\pi)^4}}
\newcommand{\mn}{{\mu\nu}}
\newcommand{\tr}{{\rm tr}}
 \newcommand{\SP}{\langle \mid S^2 \mid\rangle}
\newcommand{\Tr}{{\rm Tr}}
\newcommand{\T} {\mbox{T}}
\newcommand{\braket}[2]{\langle #1|#2\rangle}
\newcommand{\Lb}{\left(}
\newcommand{\Rb}{\right)}
\setcounter{secnumdepth}{7}
\setcounter{tocdepth}{7}
\renewcommand{\v}[1]{ \ensuremath{ {\bm{#1}} }}                  
\parskip=\itemsep               %?


%---layout fuer eine dina4 seite-------------------
\setlength{\textheight}{23cm}
\setlength{\textwidth}{178mm}
\setlength{\topmargin}{-2.5cm}
\setlength{\oddsidemargin}{1.3cm}

%\input psfig
%%%%%%%%%%%%%%%%%%%%%%%%%%%%%%%%%%%%%%%%%%%%%%%%%%%%%%%%%%%%%%%
%\renewcommand{\thefootnote}{\fnsymbol{footnote}}
%
\newcommand{\m}{\margi }
                 
\newcommand{\lash}[1]{\not\! #1 \,}
\newcommand{\nn}{\nonumber}
\newcommand{\D}{\partial}
%239
\newcommand{\h}{\frac{1}{2}}
\newcommand{\g}{{\rm g}}
\newcommand{\x}{\vec{x}}

\newcommand{\z}{\vec{z}}
\newcommand{\kv}{\vec{k}}
\newcommand{\qv}{\vec{q}}
\newcommand{\lv}{\vec{l}}
\newcommand{\mv}{\vec{m}}
%\newcommand{\q_0}{\vec{q_0}}
%\newcommand{\q'}{\vec{q}'}
%\newcommand{\v \rho}{\vec\rho}

\newcommand{\y}{{\cal Y}}
\newcommand{\el}{{\cal L}}
\newcommand{\A}{{\cal A}}
\newcommand{\Ka}{{\cal K}}
\newcommand{\al}{\alpha}
\newcommand{\be}{\beta}
\newcommand{\ep}{\varepsilon}
\newcommand{\ga}{\gamma}
\newcommand{\de}{\delta}
\newcommand{\De}{\Delta}
\newcommand{\et}{\eta}
\newcommand{\ka}{\vec{\kappa}}
\newcommand{\la}{\lambda}
\newcommand{\ph}{\varphi}
\newcommand{\ro}{\varrho}
\newcommand{\Ga}{\Gamma}
\newcommand{\om}{\omega}
\newcommand{\La}{\Lambda}
\newcommand{\tG}{\tilde{G}}

\newcommand{\pom}{I\!\!P}
\renewcommand{\theequation}{\thesection.\arabic{equation}}

%%%%%%%%%%%%%%%%%%%%%%%%%%%%%%%%%%%%%%%%%%%%%%%%%%%%%%%%
% ABBREVIATED JOURNAL NAMES
%
\def\ap#1#2#3{     {\it Ann. Phys. (NY) }{\bf #1} (19#2) #3}
\def\arnps#1#2#3{  {\it Ann. Rev. Nucl. Part. Sci. }{\bf #1} (19#2) #3}
\def\npb#1#2#3{    {\it Nucl. Phys. }{\bf B#1} (19#2) #3}
\def\plb#1#2#3{    {\it Phys. Lett. }{\bf B#1} (19#2) #3}
\def\prd#1#2#3{    {\it Phys. Rev. }{\bf D#1} (19#2) #3}
\def\prep#1#2#3{   {\it Phys. Rr ep. }{\bf #1} (19#2) #3}
\def\prl#1#2#3{    {\it Phys. Rev. Lett. }{\bf #1} (19#2) #3}
\def\ptp#1#2#3{    {\it Prog. Theor. Phys. }{\bf #1} (19#2) #3}
\def\rmp#1#2#3{    {\it Rev. Mod. Phys. }{\bf #1} (19#2) #3}
\def\zpc#1#2#3{    {\it Z. Phys. }{\bf C#1} (19#2) #3}
\def\mpla#1#2#3{   {\it Mod. Phys. Lett. }{\bf A#1} (19#2) #3}
\def\nc#1#2#3{     {\it Nuovo Cim. }{\bf #1} (19#2) #3}
\def\yf#1#2#3{     {\it Yad. Fiz. }{\bf #1} (19#2) #3}
\def\cpc#1#2#3{    {\it Comp. Phys. Commun. }{\bf #1} (19#2) #3}
\def\dis#1#2{      {\it Dissertation, }{\sf #1 } 19#2}
\def\ib#1#2#3{     {\it ibid. }{\bf #1} (19#2) #3}
\def\jpg#1#2#3{        {\it J. Phys}. {\bf G#1}#2#3}
%

%%%%%%%%%%%%%%%%%%%%%%%%%%%%%%%%%%%%%%%%%%%%%%%%%%%%%%%
%r 
%\renewcommand{\thefigure}{{\protect\bf\arabic{figure}}}
\def\thefootnote{\fnsymbol{footnote}}
\def\pom{{I\!\!P}}
\def\reg{{I\!\!R}}
\vskip1cm
%\begin{document}
%%%%%%%%%%%%%%%%%%%%%%%%%%%%%%%%%%%%%%%%%%%%%%%%%%%%%%%%%%%%%%%%%%%%%%%

\title{Renormalization group, Stochastic processes and machine learning}

\author[a]{Haowu Duan,}
\affiliation[a]{Physics Department, University of Connecticut, 2152 Hillside Road, Storrs, CT 06269, USA}

\abstract{Notes}


%\end{abstract}

\keywords{}
%PACS numbers: 12.38.-t,24.85.+p,25.75.-q
\begin{document}
\maketitle

\pagestyle{empty}


\mbox{}


\pagestyle{plain}

\setcounter{page}{1}
\author{}

\abstract{ }
\keywords{}
\dedicated{}
\preprint{}





\date{\today}

\newpage
\section{Introducing key concepts}
\subsection{Stochastic process, statistical field theory}
Here we review key results that we need for later sections in Stochastic processes and field theory.



\newpage
\subsection{Wegner-Morris equation, Polchinski's equation and gradient flow and all that}
Given a Euclidean action $S_E(\phi)$ in d dimension, we have the partition function,
\begin{equation}
\begin{split}
Z=&\int D\phi \exp\{-S_E(\phi)\}
\end{split}
\end{equation}
In practice, we don't want to deal with the full theory but effective theory at some scale $\Lambda$, where all the modes with energy greater than $\Lambda$ has been integrated out. The full theory would be defined at some initial scale $\Lambda_0$, namely, $S_E(\phi) \equiv S_{\Lambda_0}(\phi)$. Let us denote the field above some scale as $\phi_\Lambda^+$, and below that scale as $\phi_\Lambda^-$. By definition,
\begin{equation}
\begin{split}
\exp\{-S_{\Lambda}(\phi_\Lambda^-)\}
=\int D\phi_{\Lambda}^+ \exp\{-S_{\Lambda_0}(\phi)\}
\end{split}
\end{equation}
Regardless how this was achieved in practice, the partition function is independent of arbitrary scale $\Lambda$.
\begin{equation}
\begin{split}
Z=\int D\phi \exp\{-S_E(\phi)\} =\int D\phi_{\Lambda}^- \exp\{-S_{\Lambda}(\phi_\Lambda^-)\}
\end{split}
\end{equation}
Therefore, the following identity must stay true at all times,
\begin{equation}
\begin{split}
-\Lambda \frac{d}{d \Lambda}Z = \frac{d}{dt} Z =0
\end{split}
\end{equation}
where $-\Lambda \frac{d}{d \Lambda} = \frac{d}{dt }$ is the RG time, with $t=\ln(\Lambda_0/\Lambda)$. We can imagine that a class of transformation will satisfy this condition,
\begin{equation}
\frac{d}{dt }\exp\{-S_{t}\}=\int d^dx \;\frac{\delta}{\delta \phi(x)}\Big(  \Psi(t, \phi(x)) \exp\{-S_{t}\} \Big)
\end{equation}
This is known as \textit{Wegner-Morris equation}. Where $\Psi(t, \phi(x))$ is determined by RG scheme. This equation implies that under RG transformation, field variable has been reparameterized, 
\begin{equation}
\delta \phi(x)=\frac{\delta \Lambda}{ \Lambda} \Psi(t, \phi(x))
\end{equation}
{\color{red} Think about how this is true.} What gradient flow does is the following\cite{Abe:2018zdc},
\begin{equation}
\begin{split}
\partial_t \phi_t(x)=-\frac{\delta S_t}{\delta \phi(x)} ;\;\;\; \phi_{t=0}(x)=\phi_0(x) \label{Eq:GF1}
\end{split}
\end{equation}
where we update the field variable at each step based on the functional derivative of the effective action with respect to the field variable. This is also known as model A in the literature of critical dynamics. Because this is a deterministic evolution, we can denote effective field as a function of initial condition as $\phi_t[\phi_0(x)](x)$. Due to the invariance of the partition function, we can write the effective action at time t as,
\begin{equation}
\begin{split}
e^{-S_t(\phi)}=\int D\phi_0 \delta(\phi-\phi_t(\phi_0)) e^{-S_0(\phi_0)}\label{Eq:GF2}
\end{split}
\end{equation}
{\color{red} It is easy to show the invariance of partition function,
\begin{equation}
\begin{split}
&
\int D \phi e^{-S_t(\phi)}\\
=&\int D \phi \int D\phi_0 \delta(\phi-\phi_t(\phi_0)) e^{-S_0(\phi_0)}\\
=& \int D \phi_0 e^{-S_0(\phi_0)}
\end{split}
\end{equation}
}
Taking derivative on both side of the equation,
\begin{equation}
\begin{split}
\partial_t e^{-S_t(\phi)}=\int D\phi_0 \partial_t\delta(\phi-\phi_t(\phi_0)) e^{-S_0(\phi_0)}
\end{split}
\end{equation}
we can write the $\partial_t$ in terms of functional derivative using the chain rule,
\begin{equation}
\partial_t= \int dx \;\Big(\partial_t \phi(x) \Big)\frac{\delta}{\delta \phi(x)}
\end{equation}
as a result,
\begin{equation}
\begin{split}
&\partial_t e^{-S_t(\phi)}\\
=&\int D\phi_0  \int dx \;\Big(\partial_t \phi(x) \Big)\frac{\delta}{\delta \phi(x)} \delta(\phi-\phi_t(\phi_0)) e^{-S_0(\phi_0)} \\
=&\int D\phi_0  \int dx \;\Big( -\frac{\delta S_t}{\delta \phi(x)} \Big)\frac{\delta}{\delta \phi(x)} \delta(\phi-\phi_t(\phi_0)) e^{-S_0(\phi_0)}  \\
=&   \int dx \;\frac{\delta}{\delta \phi(x)}  \frac{\delta S_t}{\delta \phi(x)}\Big( \int D\phi_0   \delta(\phi-\phi_t(\phi_0)) e^{-S_0(\phi_0)} \Big) \\
=&  \int dx \;\frac{\delta}{\delta \phi(x)}  \frac{\delta S_t}{\delta \phi(x)} e^{-S_t} \label{Eq:consistency1}
\end{split}
\end{equation}
and we get,
\begin{equation}
\begin{split}
\partial_t e^{-S_t(\phi)}=\int dx \; \Big[\frac{\delta^2 S_t}{\delta \phi^2(x)}-  (\frac{\delta S_t}{\delta \phi(x)} )^2 \Big]e^{-S_t}
\end{split}
\end{equation}
expand $e^{-S_t}$, we get an equation for the effective action,
\begin{equation}
\begin{split}
\partial_t S_t(\phi)=-\int dx \; \Big[\frac{\delta^2 S_t}{\delta \phi^2(x)}-  (\frac{\delta S_t}{\delta \phi(x)} )^2 \Big]
\end{split}
\end{equation}
It is obvious that we have a UV divergence in $\frac{\delta^2 S_t}{\delta \phi^2(x)}$, which has to be regularized. 
\begin{equation}
\begin{split}
\partial_t S_t(\phi)=-\int dxdy \; \mathcal{K}_t(x,y)\Big[\frac{\delta^2 S_t}{\delta \phi(x)\delta \phi(y)}-  \frac{\delta S_t}{\delta \phi(x)} \frac{\delta S_t}{\delta \phi(y)} \Big]
\end{split}
\end{equation}
$\mathcal{K}_t(x,y)$ is a smearing kernel that remove UV modes. This is not to be confused with the kernel in Pochiski's equation, which was usually denoted as $K_t(p)$, but they are deeply related. This is because when there is $t$ dependence in the regularization kernel, there will be derivative acting on it. Their connection will be established exactly later. Suppose we separate the kinetic and interaction contribution of the action.
\begin{equation}
\begin{split}
S_t=& \int dxdy \;\; \frac{1}{2}\phi_t(x) (-\partial_x^2 +m^2 )K_t(x,y) \phi_t(y)  + I_{t} \\
      =& \int_k \frac{1}{2}\phi_t(k) (k^2 +m^2)K_t(k)\phi_t(-k) + I_{t} 
\end{split}
\end{equation}
when unregularized, $K_t(x,y) =\Rigtharrow \delta(x-y)$, and the derivative formally act on $\delta(x-y)$. So far we did not specify the functional form of the smearing kernel. But we know that it has to remove high energy modes. Recall the relation between RG time and the momentum scale $\Lambda$ with initial condition $\Lambda_0$, $t=\ln \frac{\Lambda_0}{\Lambda}$. If $K_t(k)$ is a sharp cut-off, we can pick $K_t(k)=\theta(\Lambda^2_0e^{-2t}-k^2)$, for a smoother choice, we can use $K_t(k)=\exp(-k^2/\Lambda^2_0e^{-2t})$, which corresponds to a gaussian coarse graining procedure. Now we go through the similar procedure as Eq.~\eqref{Eq:consistency1} to derive the equation.
\begin{equation}
\begin{split}
&\partial_t e^{-S_t(\phi)}\\
=&\int_k \Big[\frac{1}{2}\phi_t(k) (k^2 +m^2)\phi_t(-k) (\partial_t K_t(k))+ \partial_t I_{t}  \Big]e^{-S_t(\phi)}
 \\
\label{Eq:consistency2}
\end{split}
\end{equation}
Now the task is, without using gradient flow condition Eq.~\ref{Eq:GF2}, how do we make sure the invariance of the partition function. 
\begin{equation}
\begin{split}
\partial_t Z 
=& \int D\phi_t  \;\partial_t e^{-S_t(\phi_t)} \\
=& \int D\phi_t \int_k \Big[\frac{1}{2}\phi_t(k) (k^2 +m^2)\phi_t(-k) \partial_t K_t(k)+ \partial_t I_{t}  \Big]e^{-S_t(\phi)}\\
=& 0
\end{split}
\end{equation}
This means, the interaction part of the action must evolves according to,
\begin{equation}
\langle \partial_t I_{t} \rangle=-\int_k \Big[\frac{1}{2} (k^2 +m^2)(\partial_t K_t(k)) \langle \phi_t(k) \phi_t(-k)  \rangle
\end{equation}
Now we have to think about what to do with $\langle \phi_t(k) \phi_t(-k)  \rangle$, and if we can somehow express it in term of $S_{t; int}$ alone, we will get an equation for $S_{t; int}$. 
%By definition, the two point function is,
%\begin{equation}
%\begin{split}
 %\langle \phi_t(k) \phi_t(-k)  \rangle=\frac{1}{Z} \frac{\delta^2 Z(J)}{\delta J(-k)\delta J(k)}\Big|_{J=0}
%\end{split}
%\end{equation}
We note that $K_t(k)=\exp(-k^2/\Lambda^2_0e^{-2t})$,
\begin{equation}
\partial_t K_t(k)=(-2k^2/\Lambda^2_0e^{-2t})   \exp(-k^2/\Lambda^2_0e^{-2t})  
\end{equation}
when $k^2\ll \Lambda^2_0e^{-2t}$, the derivative is almost zero, while for $k^2\gg \Lambda^2_0e^{-2t}$, there is strong exp suppression. As a result, only $k^2/\Lambda^2_0e^{-2t} \sim \mathcal{O}(1)$ will contribute, we can think of this as smeared $\delta(k^2/\Lambda^2_0e^{-2t} -1)$. For convenience, we let $\Lambda_0e^{-t}=\Lambda$. To extract the effect due to the integration over the momentum shell around $\Lambda$, we first reduce the scale by $\Delta$ then take the continue limit. Let all the field modes to be separated into $\phi_\Lambda=\phi_{l}+ \phi_{h}$, where $h$ stands for high energy modes to be integrated out.
\begin{equation}
\begin{split}
&I_{\Lambda}(\phi_\Lambda)\\
=&I_\Lambda(\phi_{l}+ \phi_{h})\\
=&I_\Lambda(\phi_{l})+\int_{k\in shell} \frac{\delta I_\Lambda}{\delta \phi_{h}(k)}\Big|_{\phi_{h}(k)=0} \phi_{h}(k)\\
+ & \frac{1}{2}\int_{k\in shell} \frac{\delta^2 I_\Lambda}{\delta \phi_{h}(k)\delta \phi_{h}(-k)}\Big|_{\phi_{h}(k)=0} \phi_{h}(k)\phi_{h}(-k)
\end{split}
\end{equation}
we take,
\begin{equation}
\begin{split}
 &\langle \phi_t(k) \phi_t(-k)  \rangle_{k\in shell} \rightarrow \int D\phi_{h}\; \phi_{h}(k) \phi_h(-k) e^{-S_{shell}}
\end{split}
\end{equation}
where,
\begin{equation}
\begin{split}
S_{shell}=& \int_{k\in shell} \frac{1}{2}\phi_\Lambda(k) (k^2 +m^2)K_\Lambda(k)\phi_\Lambda(-k) + \int_{k\in shell} \frac{\delta I_\Lambda}{\delta \phi_{h}(k)}\Big|_{\phi_{h}(k)=0} \phi_{h}(k)\\
+ &\frac{1}{2}\int_{k\in shell} \frac{\delta^2 I_\Lambda}{\delta \phi_{h}(k)\delta \phi_{h}(-k)}\Big|_{\phi_{h}(k)=0} \phi_{h}(k)\phi_{h}(-k)\\ \\ \\
\end{split}
\end{equation}
{\color{red} We have to be very clear about what kind of correction was included in the expectation value. This is highly non-trivial}


\subsection{An exercise to internalize the concepts: zero dimensional field theory}
In zero dimension, field variable $\phi(x)$ reduces to random variable $\phi$. We define the probability density $P_t(\phi)$,
\begin{equation}
\begin{split}
P_t(\phi)=\frac{e^{-E_t(\phi)}}{Z_0}
\end{split}
\end{equation}
where t is again the RG time. And the partition function is RG invariant. We are interested in how $P_t(\phi)$ or $\rho_t(\phi)=e^{-E_t(\phi)}$ evolves with time. It is non-trivial to write down the exact equation from first principle, but nevertheless, one can write to write down the equation that works. {\color{red}For now, let me just put into the equations}. Suppose the evolution of $\rho_t(\phi)$ satisfy the following Fokker-Planck equation of one-dimensional Ornstein-Uhlenbeck process. 
\begin{equation}
\begin{split}
\frac{\partial \rho(\phi)}{\partial t}&=\frac{\partial }{\partial \phi} (\omega \phi+ D\frac{\partial }{\partial \phi}) \rho_t(\phi) \\
&=\frac{\partial }{\partial \phi} \omega \phi \rho_t(\phi) +D \frac{\partial^2 }{\partial \phi^2} \rho_t(\phi) \label{Eq:WK}
\end{split}
\end{equation}
recall that in general FP is of the form,
\begin{equation}
\begin{split}
\frac{\partial \rho(\phi)}{\partial t}=\frac{\partial }{\partial \phi} A(\phi, t) \rho_t(\phi) +\frac{1}{2}\frac{\partial^2 }{\partial \phi^2} B(\phi,t) \rho_t(\phi)
\end{split}
\end{equation}
This corresponds to a specific choice of drift and diffusion coefficients. Eq.~\eqref{Eq:WK} has the following solution,
\begin{equation}
\begin{split}
\rho_t(\phi)=\int d\phi\; P(\phi, t | \phi_0,0) \rho_0(\phi_0)
\end{split}
\end{equation}
{\color{red} The fact that WK only wrote the equation for un-normalized probability means their FK equation has to be linear. The question is then whether non-linear FK has any relevance for applications in field theory and Machine learning.} Where the propagator is given by,
\begin{equation}
\begin{split}
P(\phi, t | \phi_0,0)=\Big[\frac{\omega}{2\pi D(1-e^{-2\omega t})} \Big]^{1/2} \exp\{ -\frac{\omega(\phi-e^{-\omega t}\phi_0)^2}{2D(1-e^{-2\omega t})}\}
\end{split}
\end{equation}
One can clearly see at $t=0$ limit, we get $\lim_{t\rightarrow 0}\rho_t(\phi)=\rho_0(\phi_0)$ because the propagator collapsed into $\delta(\phi-\phi_0)$. On the other hand, $\rho_{\infty}(\phi)$ is a Gaussian distribution without any memories of the initial condition. The WK RG scheme can be generalized to higher dimension but we will discuss this later along Carosso RG scheme. It is easy to check that,
\begin{equation}
\int d\phi\; \rho_t(\phi)=\int d\phi_0\int d\phi\; P(\phi, t | \phi_0,0)\rho_0(\phi_0)=\int d\phi_0\; \rho_0(\phi_0)=Z_0
\end{equation}
Therefore, partition function is RG invariant. What Wegner-Morris proposed is that RG corresponds to a class of reparameterization transformation,
\begin{equation}
\begin{split}
\phi'=\phi+\Psi_t(\phi) dt \;\;\; \Rightarrow \;\;\; \frac{\delta \phi'}{\delta \phi}=1+ \frac{\delta \Psi_t(\phi)}{\delta \phi} dt
\end{split}
\end{equation}
we plug this transformation into the evolution of $Z_t$, and demand its invariance.
\begin{equation}
\begin{split}
\int d\phi' P_{t+d t}(\phi') - \int d\phi P_{t}(\phi) =0
\end{split}
\end{equation}
we than extract the part that is linear in $dt$. 
\begin{equation}
\begin{split}
\rho_{t+d t}(\phi') &=\rho_{t}(\phi)+ dt\Big[ \partial_t\rho_{t}(\phi)+\Psi_t(\phi) \partial_{\phi}\rho_{t}(\phi) \Big] \\
d\phi'&=d\phi+d\Psi_t(\phi) dt =d\phi+\partial_{\phi}\Psi_t(\phi) dtd \phi
\end{split}
\end{equation}
Therefore,
\begin{equation}
\begin{split}
&\int d\phi \; \Big(\partial_t\rho_{t}(\phi)+\Psi_t(\phi) \partial_{\phi}\rho_{t}(\phi) +\rho_{t}(\phi) \partial_{\phi}\Psi_t(\phi) \Big) \\
=&\int d\phi \; \Big(\partial_t\rho_{t}(\phi)+ \partial_{\phi}\Psi_t(\phi)\rho_{t}(\phi) \Big) \\
=& 0
\end{split}
\end{equation}
Therefore we have $\partial_tP_{t}(\phi)+ \partial_{\phi}\Psi_t(\phi)P_{t}(\phi)=0$. However, one might wonder this is not the most general case as the result of this equation can be a total derivative of any function with vanishing boundary condition. However, since our choice of $\Psi_t(\phi)$ is not fix and it represent a class of transformation. The total derivative can be absorbed into the definition of $\Psi_t$. One can generalize from zero dimension to d dimension and obtain the Wegner-Morris equation. It's clear that FK of WK is a special case of Wegner-Morris equation,
\begin{equation}
\begin{split}
\Psi_t(\phi)=\omega \phi+ D\frac{\partial \ln \rho_t(\phi)}{\partial \phi}
\end{split}
\end{equation}
There are claims that the irreversibility of RG can be made sure if $\Psi_t(\phi,\rho_t)$ has specific $\rho_t$ dependence, namely the distribution at previous time. {\color{red} However, even for such types of equation, one can still derive time-reversed stochastic differential equation. So I am not so sure about this part. Also even if one add the $\rho_t$, there will be multiple discretization scheme that violates the time irreversibility but still have the same continue limit. } Next, we discuss Carosso RG and demonstrate that it has a clear interpretation of coarse graining in coordinate space. Let us start with the the Langevin equation associated with the Ornstein-Uhlenbeck process,
\begin{equation}
\begin{split}
\partial_t \phi_t=-\omega \phi_t +\xi_t
\end{split}
\end{equation}
where the $\xi_t$ satisfy,
\begin{equation}
\begin{split}
\langle \xi_t \rangle=0 \;\; ;\langle \xi_t \xi_{t'} \rangle=2D\delta(t-t')
\end{split}
\end{equation}
The solution of the Langevin equation is,
\begin{equation}
\begin{split}
\phi_t=e^{-\omega t}\phi_0+ \int_0^t d\tau e^{-(t-\tau)\omega} \xi_{\tau}
\end{split}
\end{equation}
Carrosso RG corresponds to the following \textit{Stochastic heat equation}, 
\begin{equation}
\partial_t \phi_t(x)=\partial^2_\mu \phi_t(x) +\xi_t(x)
\end{equation}
where the noises are normalized as,
\begin{equation}
\langle \xi_t(x) \xi_{t'}(y) \rangle=2D \delta^{d}(x-y)\delta(t-t')
\end{equation}
This implies that in momentum space, $\omega=k^2$, where k is the momentum of the modes. For simplicity, let's donote $\partial^2_\mu =\Delta$. Operatorially, the solution is,
\begin{equation}
\phi_t(x)=\int d^dy \langle  x|e^{t\Delta}| y\rangle \phi_0(y)+ \int_0^t d d\tau d^d y \langle  x|e^{(t-\tau)\Delta}| y\rangle \xi_\tau(y)
\end{equation}
where the heat kernel is well-known result,
\begin{equation}
\begin{split}
 \langle  x|e^{t\Delta}| y\rangle=\frac{1}{[4\pi t]^{d/2}} e^{-\frac{(x-y)^2}{4t}}=\int \frac{d ^dk}{(2\pi)^d} e^{-i(x-y)\cdot k}e^{-k^2 t}
\end{split}
\end{equation}
and we can immediately identify the real-space \textit{Coarse Graining} with the Gaussian smear $e^{-\frac{(x-y)^2}{4t}}$. Given the solution of the langevin equation, we can construct the corresponding Fokker-Planck equation for the probability density. Before we do that, let's further clarify the noise, namely,


\subsection{Gradient flow and renormalization group}

\subsection{Gaussian process and its connection to attention mechanism}

\subsection{Relavent ML architectures}


\newpage
\section{Renormalization as an optimization problem}


\newpage
\section{Renormalizing diffusion model and multi-scale diffusion model}

\bibliographystyle{JHEP}
\bibliography{Notes}

\end{document}
